\section{IBM graph-construction and scale results}
\label{sec:s-ibm-graph-scale}

\subsection{Dependence-aware matched design}

Each candidate is paired to SAGE-edge h64 by variant, protocol, and seed. The four fixed HI/LI-by-protocol differences are averaged within seed before bootstrap and Wilcoxon inference, so \(n=10\), not 40. The following tables report every candidate-size aggregate for AUPRC, AUROC, and F1; all 208 context-specific sensitivities remain machine-readable.

\subsection{AUPRC effects}

\begingroup
\footnotesize
\setlength{\tabcolsep}{3pt}
\renewcommand{\arraystretch}{1.10}
\begin{longtable}{@{}p{0.08\textwidth} p{0.18\textwidth} r r p{0.24\textwidth} r r p{0.10\textwidth}@{}}
\caption{Dependence-aware matched IBM AUPRC effects relative to SAGE-edge h64.}\label{tab:s-ibm-effect-auprc}\\
\toprule
\textbf{Size} & \textbf{Candidate} & \textbf{Reference} & \textbf{Candidate} & \textbf{$\Delta$ [95\% CI]} & \textbf{$d_z$} & \textbf{Holm p} & \textbf{Context signs} \\
\midrule
\endfirsthead
\caption[]{Dependence-aware matched IBM AUPRC effects relative to SAGE-edge h64. (continued)}\\
\toprule
\textbf{Size} & \textbf{Candidate} & \textbf{Reference} & \textbf{Candidate} & \textbf{$\Delta$ [95\% CI]} & \textbf{$d_z$} & \textbf{Holm p} & \textbf{Context signs} \\
\midrule
\endhead
\midrule
\multicolumn{8}{r}{\footnotesize Continued on next page}\\
\endfoot
\bottomrule
\endlastfoot
Medium & DegreeCap & 0.02844 & 0.00758 & -0.02086 [-0.02248, -0.01935] & -7.854 & 0.0117 & +0 / -4 / 0:0 \\
Medium & DegreeOnly & 0.02844 & 0.01783 & -0.01061 [-0.01287, -0.00821] & -2.674 & 0.0117 & +0 / -4 / 0:0 \\
Medium & NoEdge & 0.02844 & 0.01638 & -0.01206 [-0.01395, -0.01022] & -3.783 & 0.0117 & +0 / -4 / 0:0 \\
Medium & RecentWindow & 0.02844 & 0.01520 & -0.01324 [-0.01543, -0.01092] & -3.408 & 0.0117 & +0 / -4 / 0:0 \\
Medium & Sender-receiver alias & 0.02844 & 0.02844 & +0.00000 [0.00000, 0.00000] & 0.621 & 1.0000 & +0 / -0 / 0:4 \\
Medium & ShuffledEdge & 0.02844 & 0.01636 & -0.01208 [-0.01324, -0.01089] & -6.042 & 0.0117 & +0 / -4 / 0:0 \\
Small & DegreeCap & 0.00752 & 0.01182 & +0.00430 [-0.00036, 0.00971] & 0.504 & 0.4648 & +3 / -1 / 0:0 \\
Small & DegreeOnly & 0.00752 & 0.00381 & -0.00371 [-0.00547, -0.00202] & -1.264 & 0.0156 & +0 / -4 / 0:0 \\
Small & GINE h64 & 0.00752 & 0.01942 & +0.01190 [0.00934, 0.01447] & 2.697 & 0.0137 & +4 / -0 / 0:0 \\
Small & NoEdge & 0.00752 & 0.00485 & -0.00267 [-0.00423, -0.00109] & -0.973 & 0.0820 & +0 / -4 / 0:0 \\
Small & RecentWindow & 0.00752 & 0.00418 & -0.00334 [-0.00553, -0.00148] & -0.976 & 0.0137 & +0 / -4 / 0:0 \\
Small & Sender-receiver alias & 0.00752 & 0.00752 & +0.00000 [0.00000, 0.00000] & 0.000 & 1.0000 & +0 / -0 / 0:4 \\
Small & ShuffledEdge & 0.00752 & 0.00422 & -0.00330 [-0.00487, -0.00197] & -1.329 & 0.0137 & +0 / -4 / 0:0 \\
\end{longtable}
\noindent\footnotesize\emph{Scope and reading.} Inference uses ten seed blocks after averaging four fixed HI/LI-by-protocol contexts within seed. The context-sign column summarizes heterogeneity; the complete 208 context rows remain in the artifact. The sender-receiver row is an identity diagnostic.\par
\endgroup


NoEdge changes mean AUPRC by \(-0.00267\) on Small and \(-0.01206\) on Medium. All four context means are negative at each scale, but only Medium survives Holm correction (\(p=0.0117\), versus 0.0820 on Small). ShuffledEdge changes AUPRC by \(-0.00330\) and \(-0.01208\), with Holm values 0.0137 and 0.0117. These matched controls support a bounded role for transaction-feature alignment in this implementation, not a causal statement about all AML graphs.

GINE improves Small AUPRC by 0.01190 (95\% CI \([0.00934,0.01447]\), Holm \(p=0.0137\)) with all four context means positive. Medium GINE is absent because it exhausted T4 memory; Small-only performance is not averaged into a complete Small/Medium claim.

\subsection{AUROC effects}

\begingroup
\footnotesize
\setlength{\tabcolsep}{3pt}
\renewcommand{\arraystretch}{1.10}
\begin{longtable}{@{}p{0.08\textwidth} p{0.18\textwidth} r r p{0.24\textwidth} r r p{0.10\textwidth}@{}}
\caption{Dependence-aware matched IBM AUROC effects relative to SAGE-edge h64.}\label{tab:s-ibm-effect-auroc}\\
\toprule
\textbf{Size} & \textbf{Candidate} & \textbf{Reference} & \textbf{Candidate} & \textbf{$\Delta$ [95\% CI]} & \textbf{$d_z$} & \textbf{Holm p} & \textbf{Context signs} \\
\midrule
\endfirsthead
\caption[]{Dependence-aware matched IBM AUROC effects relative to SAGE-edge h64. (continued)}\\
\toprule
\textbf{Size} & \textbf{Candidate} & \textbf{Reference} & \textbf{Candidate} & \textbf{$\Delta$ [95\% CI]} & \textbf{$d_z$} & \textbf{Holm p} & \textbf{Context signs} \\
\midrule
\endhead
\midrule
\multicolumn{8}{r}{\footnotesize Continued on next page}\\
\endfoot
\bottomrule
\endlastfoot
Medium & DegreeCap & 0.89599 & 0.87608 & -0.01991 [-0.03384, -0.00674] & -0.867 & 0.0742 & +1 / -3 / 0:0 \\
Medium & DegreeOnly & 0.89599 & 0.80754 & -0.08845 [-0.09960, -0.07698] & -4.644 & 0.0117 & +0 / -4 / 0:0 \\
Medium & NoEdge & 0.89599 & 0.80501 & -0.09098 [-0.09994, -0.08143] & -5.730 & 0.0117 & +0 / -4 / 0:0 \\
Medium & RecentWindow & 0.89599 & 0.92124 & +0.02525 [0.01610, 0.03484] & 1.576 & 0.0117 & +4 / -0 / 0:0 \\
Medium & Sender-receiver alias & 0.89599 & 0.89599 & +0.00000 [0.00000, 0.00000] & 0.089 & 1.0000 & +0 / -0 / 0:4 \\
Medium & ShuffledEdge & 0.89599 & 0.79679 & -0.09920 [-0.11548, -0.08282] & -3.543 & 0.0117 & +0 / -4 / 0:0 \\
Small & DegreeCap & 0.84528 & 0.86881 & +0.02354 [0.00573, 0.04076] & 0.788 & 0.1465 & +3 / -1 / 0:0 \\
Small & DegreeOnly & 0.84528 & 0.73127 & -0.11401 [-0.13059, -0.09730] & -4.039 & 0.0137 & +0 / -4 / 0:0 \\
Small & GINE h64 & 0.84528 & 0.85618 & +0.01090 [-0.01596, 0.03856] & 0.231 & 1.0000 & +2 / -2 / 0:0 \\
Small & NoEdge & 0.84528 & 0.73569 & -0.10959 [-0.12200, -0.09520] & -4.844 & 0.0137 & +0 / -4 / 0:0 \\
Small & RecentWindow & 0.84528 & 0.79306 & -0.05222 [-0.06542, -0.03977] & -2.426 & 0.0137 & +0 / -4 / 0:0 \\
Small & Sender-receiver alias & 0.84528 & 0.84528 & +0.00000 [0.00000, 0.00000] & 0.000 & 1.0000 & +0 / -0 / 0:4 \\
Small & ShuffledEdge & 0.84528 & 0.72068 & -0.12460 [-0.13473, -0.11382] & -6.880 & 0.0137 & +0 / -4 / 0:0 \\
\end{longtable}
\noindent\footnotesize\emph{Scope and reading.} Inference uses ten seed blocks after averaging four fixed HI/LI-by-protocol contexts within seed. The context-sign column summarizes heterogeneity; the complete 208 context rows remain in the artifact. The sender-receiver row is an identity diagnostic.\par
\endgroup


AUROC sometimes moves in the opposite direction from AUPRC. RecentWindow on Medium lowers AUPRC but raises mean AUROC by 0.02525. This tradeoff is retained rather than reduced to a winner label.

\subsection{F1 effects}

\begingroup
\footnotesize
\setlength{\tabcolsep}{3pt}
\renewcommand{\arraystretch}{1.10}
\begin{longtable}{@{}p{0.08\textwidth} p{0.18\textwidth} r r p{0.24\textwidth} r r p{0.10\textwidth}@{}}
\caption{Dependence-aware matched IBM F1 effects relative to SAGE-edge h64.}\label{tab:s-ibm-effect-f1}\\
\toprule
\textbf{Size} & \textbf{Candidate} & \textbf{Reference} & \textbf{Candidate} & \textbf{$\Delta$ [95\% CI]} & \textbf{$d_z$} & \textbf{Holm p} & \textbf{Context signs} \\
\midrule
\endfirsthead
\caption[]{Dependence-aware matched IBM F1 effects relative to SAGE-edge h64. (continued)}\\
\toprule
\textbf{Size} & \textbf{Candidate} & \textbf{Reference} & \textbf{Candidate} & \textbf{$\Delta$ [95\% CI]} & \textbf{$d_z$} & \textbf{Holm p} & \textbf{Context signs} \\
\midrule
\endhead
\midrule
\multicolumn{8}{r}{\footnotesize Continued on next page}\\
\endfoot
\bottomrule
\endlastfoot
Medium & DegreeCap & 0.00809 & 0.01396 & +0.00587 [0.00317, 0.00841] & 1.308 & 0.0352 & +4 / -0 / 0:0 \\
Medium & DegreeOnly & 0.00809 & 0.00834 & +0.00026 [-0.00292, 0.00282] & 0.052 & 1.0000 & +1 / -3 / 0:0 \\
Medium & NoEdge & 0.00809 & 0.00739 & -0.00070 [-0.00339, 0.00135] & -0.170 & 1.0000 & +1 / -3 / 0:0 \\
Medium & RecentWindow & 0.00809 & 0.01228 & +0.00419 [0.00242, 0.00583] & 1.416 & 0.0352 & +3 / -1 / 0:0 \\
Medium & Sender-receiver alias & 0.00809 & 0.00809 & +0.00000 [0.00000, 0.00000] & 0.000 & 1.0000 & +0 / -0 / 0:4 \\
Medium & ShuffledEdge & 0.00809 & 0.00723 & -0.00086 [-0.00368, 0.00136] & -0.199 & 1.0000 & +1 / -3 / 0:0 \\
Small & DegreeCap & 0.01058 & 0.01532 & +0.00473 [0.00204, 0.00747] & 1.006 & 0.0781 & +4 / -0 / 0:0 \\
Small & DegreeOnly & 0.01058 & 0.00693 & -0.00365 [-0.00518, -0.00226] & -1.458 & 0.0137 & +0 / -4 / 0:0 \\
Small & GINE h64 & 0.01058 & 0.00852 & -0.00206 [-0.00371, -0.00052] & -0.768 & 0.2520 & +0 / -4 / 0:0 \\
Small & NoEdge & 0.01058 & 0.00729 & -0.00329 [-0.00470, -0.00201] & -1.452 & 0.0137 & +0 / -4 / 0:0 \\
Small & RecentWindow & 0.01058 & 0.00862 & -0.00197 [-0.00385, -0.00029] & -0.645 & 0.2617 & +1 / -3 / 0:0 \\
Small & Sender-receiver alias & 0.01058 & 0.01058 & +0.00000 [0.00000, 0.00000] & 0.000 & 1.0000 & +0 / -0 / 0:4 \\
Small & ShuffledEdge & 0.01058 & 0.00723 & -0.00335 [-0.00492, -0.00195] & -1.327 & 0.0195 & +0 / -4 / 0:0 \\
\end{longtable}
\noindent\footnotesize\emph{Scope and reading.} Inference uses ten seed blocks after averaging four fixed HI/LI-by-protocol contexts within seed. The context-sign column summarizes heterogeneity; the complete 208 context rows remain in the artifact. The sender-receiver row is an identity diagnostic.\par
\endgroup


GINE's Small F1 change is \(-0.00206\), but its corrected test is not significant (\(p=0.252\)). DegreeCap and RecentWindow can improve fixed-threshold F1 while lowering AUPRC. These results are specific to the saved threshold of 0.5.

\begin{figure}[tbp]
    \centering
    \includegraphics[width=0.88\textwidth]{fig05_matched_ablation_effects.pdf}
    \caption{Matched IBM effects relative to SAGE-edge h64. Intervals use ten seed blocks after fixed-context aggregation; Small-only GINE and blocked Medium cells are not pooled.}
    \label{fig:s-ibm-effects-v2}
\end{figure}

\subsection{Configuration-specific runtime and feasibility}

Table~\ref{tab:s-ibm-runtime-small} gives every Small configuration its own AUPRC and elapsed runner time.

\begingroup
\footnotesize
\setlength{\tabcolsep}{3pt}
\renewcommand{\arraystretch}{1.10}
\begin{longtable}{@{}p{0.11\textwidth} p{0.13\textwidth} p{0.19\textwidth} p{0.16\textwidth} r p{0.16\textwidth}@{}}
\caption{IBM Small configuration-level AUPRC, runtime, and feasibility.}\label{tab:s-ibm-runtime-small}\\
\toprule
\textbf{Variant} & \textbf{Protocol} & \textbf{Configuration} & \textbf{Status} & \textbf{AUPRC / seconds} & \textbf{Within-cell treatment} \\
\midrule
\endfirsthead
\caption[]{IBM Small configuration-level AUPRC, runtime, and feasibility. (continued)}\\
\toprule
\textbf{Variant} & \textbf{Protocol} & \textbf{Configuration} & \textbf{Status} & \textbf{AUPRC / seconds} & \textbf{Within-cell treatment} \\
\midrule
\endhead
\midrule
\multicolumn{6}{r}{\footnotesize Continued on next page}\\
\endfoot
\bottomrule
\endlastfoot
HI SMALL & Early-to-late & DegreeCap & \texttt{MEASURED\_\allowbreak{}FULL10} & 0.01285 / 27.89 & Pareto \\
HI SMALL & Early-to-late & DegreeOnly & \texttt{MEASURED\_\allowbreak{}FULL10} & 0.00513 / 31.61 & Measured, non-Pareto \\
HI SMALL & Early-to-late & GINE h64 & \texttt{MEASURED\_\allowbreak{}FULL10} & 0.02608 / 264.47 & Pareto \\
HI SMALL & Early-to-late & NoEdge & \texttt{MEASURED\_\allowbreak{}FULL10} & 0.00387 / 35.00 & Measured, non-Pareto \\
HI SMALL & Early-to-late & RecentWindow & \texttt{MEASURED\_\allowbreak{}FULL10} & 0.00496 / 16.64 & Pareto \\
HI SMALL & Early-to-late & SAGE-edge h64 & \texttt{MEASURED\_\allowbreak{}FULL10} & 0.01025 / 33.27 & Measured, non-Pareto \\
HI SMALL & Early-to-late & ShuffledEdge & \texttt{MEASURED\_\allowbreak{}FULL10} & 0.00412 / 32.52 & Measured, non-Pareto \\
HI SMALL & Late-window & DegreeCap & \texttt{MEASURED\_\allowbreak{}FULL10} & 0.02574 / 39.06 & Pareto \\
HI SMALL & Late-window & DegreeOnly & \texttt{MEASURED\_\allowbreak{}FULL10} & 0.00624 / 44.52 & Measured, non-Pareto \\
HI SMALL & Late-window & GINE h64 & \texttt{MEASURED\_\allowbreak{}FULL10} & 0.03125 / 363.32 & Pareto \\
HI SMALL & Late-window & NoEdge & \texttt{MEASURED\_\allowbreak{}FULL10} & 0.01052 / 48.70 & Measured, non-Pareto \\
HI SMALL & Late-window & RecentWindow & \texttt{MEASURED\_\allowbreak{}FULL10} & 0.00682 / 29.69 & Pareto \\
HI SMALL & Late-window & SAGE-edge h64 & \texttt{MEASURED\_\allowbreak{}FULL10} & 0.01100 / 47.21 & Measured, non-Pareto \\
HI SMALL & Late-window & ShuffledEdge & \texttt{MEASURED\_\allowbreak{}FULL10} & 0.00752 / 50.12 & Measured, non-Pareto \\
LI SMALL & Early-to-late & DegreeCap & \texttt{MEASURED\_\allowbreak{}FULL10} & 0.00389 / 42.79 & Pareto \\
LI SMALL & Early-to-late & DegreeOnly & \texttt{MEASURED\_\allowbreak{}FULL10} & 0.00168 / 48.76 & Measured, non-Pareto \\
LI SMALL & Early-to-late & GINE h64 & \texttt{MEASURED\_\allowbreak{}FULL10} & 0.00974 / 494.76 & Pareto \\
LI SMALL & Early-to-late & NoEdge & \texttt{MEASURED\_\allowbreak{}FULL10} & 0.00217 / 46.50 & Measured, non-Pareto \\
LI SMALL & Early-to-late & RecentWindow & \texttt{MEASURED\_\allowbreak{}FULL10} & 0.00225 / 27.39 & Pareto \\
LI SMALL & Early-to-late & SAGE-edge h64 & \texttt{MEASURED\_\allowbreak{}FULL10} & 0.00315 / 50.84 & Measured, non-Pareto \\
LI SMALL & Early-to-late & ShuffledEdge & \texttt{MEASURED\_\allowbreak{}FULL10} & 0.00199 / 46.62 & Measured, non-Pareto \\
LI SMALL & Late-window & DegreeCap & \texttt{MEASURED\_\allowbreak{}FULL10} & 0.00479 / 53.44 & Pareto \\
LI SMALL & Late-window & DegreeOnly & \texttt{MEASURED\_\allowbreak{}FULL10} & 0.00218 / 59.71 & Measured, non-Pareto \\
LI SMALL & Late-window & GINE h64 & \texttt{MEASURED\_\allowbreak{}FULL10} & 0.01060 / 675.32 & Pareto \\
LI SMALL & Late-window & NoEdge & \texttt{MEASURED\_\allowbreak{}FULL10} & 0.00283 / 58.44 & Measured, non-Pareto \\
LI SMALL & Late-window & RecentWindow & \texttt{MEASURED\_\allowbreak{}FULL10} & 0.00268 / 35.34 & Pareto \\
LI SMALL & Late-window & SAGE-edge h64 & \texttt{MEASURED\_\allowbreak{}FULL10} & 0.00568 / 65.92 & Pareto \\
LI SMALL & Late-window & ShuffledEdge & \texttt{MEASURED\_\allowbreak{}FULL10} & 0.00325 / 58.06 & Measured, non-Pareto \\
\end{longtable}
\noindent\footnotesize\emph{Scope and reading.} Pareto treatment is computed only within the same variant and protocol. Runtime is runner elapsed time, not a hardware-normalized complexity or energy measure. Blocked rows remain nonnumeric.\par
\endgroup


GINE takes approximately 9.1 times the h64 reference runtime in the aggregated Small comparison. Its higher AUPRC therefore does not imply universal deployability or dominance under every resource objective.

Table~\ref{tab:s-ibm-runtime-medium} gives the corresponding Medium cells and retains blocked GINE rows as nonnumeric statuses.

\begingroup
\footnotesize
\setlength{\tabcolsep}{3pt}
\renewcommand{\arraystretch}{1.10}
\begin{longtable}{@{}p{0.11\textwidth} p{0.13\textwidth} p{0.19\textwidth} p{0.16\textwidth} r p{0.16\textwidth}@{}}
\caption{IBM Medium configuration-level AUPRC, runtime, and feasibility.}\label{tab:s-ibm-runtime-medium}\\
\toprule
\textbf{Variant} & \textbf{Protocol} & \textbf{Configuration} & \textbf{Status} & \textbf{AUPRC / seconds} & \textbf{Within-cell treatment} \\
\midrule
\endfirsthead
\caption[]{IBM Medium configuration-level AUPRC, runtime, and feasibility. (continued)}\\
\toprule
\textbf{Variant} & \textbf{Protocol} & \textbf{Configuration} & \textbf{Status} & \textbf{AUPRC / seconds} & \textbf{Within-cell treatment} \\
\midrule
\endhead
\midrule
\multicolumn{6}{r}{\footnotesize Continued on next page}\\
\endfoot
\bottomrule
\endlastfoot
HI MEDIUM & Both Planned Protocols & GINE h64 & \texttt{RESOURCE\_\allowbreak{}BLOCKED\_\allowbreak{}T4\_\allowbreak{}CUDA\_\allowbreak{}OOM} & -- & Unmeasured \\
HI MEDIUM & Early-to-late & DegreeCap & \texttt{MEASURED\_\allowbreak{}FULL10} & 0.01186 / 216.79 & Measured, non-Pareto \\
HI MEDIUM & Early-to-late & DegreeOnly & \texttt{MEASURED\_\allowbreak{}FULL10} & 0.02524 / 238.96 & Pareto \\
HI MEDIUM & Early-to-late & NoEdge & \texttt{MEASURED\_\allowbreak{}FULL10} & 0.02447 / 248.13 & Measured, non-Pareto \\
HI MEDIUM & Early-to-late & RecentWindow & \texttt{MEASURED\_\allowbreak{}FULL10} & 0.01866 / 128.15 & Pareto \\
HI MEDIUM & Early-to-late & SAGE-edge h64 & \texttt{MEASURED\_\allowbreak{}FULL10} & 0.04242 / 249.73 & Pareto \\
HI MEDIUM & Early-to-late & ShuffledEdge & \texttt{MEASURED\_\allowbreak{}FULL10} & 0.02388 / 255.57 & Measured, non-Pareto \\
HI MEDIUM & Late-window & DegreeCap & \texttt{MEASURED\_\allowbreak{}FULL10} & 0.01274 / 261.86 & Measured, non-Pareto \\
HI MEDIUM & Late-window & DegreeOnly & \texttt{MEASURED\_\allowbreak{}FULL10} & 0.03262 / 290.36 & Pareto \\
HI MEDIUM & Late-window & NoEdge & \texttt{MEASURED\_\allowbreak{}FULL10} & 0.02989 / 298.87 & Measured, non-Pareto \\
HI MEDIUM & Late-window & RecentWindow & \texttt{MEASURED\_\allowbreak{}FULL10} & 0.02895 / 153.10 & Pareto \\
HI MEDIUM & Late-window & SAGE-edge h64 & \texttt{MEASURED\_\allowbreak{}FULL10} & 0.04998 / 301.12 & Pareto \\
HI MEDIUM & Late-window & ShuffledEdge & \texttt{MEASURED\_\allowbreak{}FULL10} & 0.02970 / 311.00 & Measured, non-Pareto \\
LI MEDIUM & Both Planned Protocols & GINE h64 & \texttt{RESOURCE\_\allowbreak{}BLOCKED\_\allowbreak{}T4\_\allowbreak{}CUDA\_\allowbreak{}OOM} & -- & Unmeasured \\
LI MEDIUM & Early-to-late & DegreeCap & \texttt{MEASURED\_\allowbreak{}FULL10} & 0.00316 / 211.08 & Measured, non-Pareto \\
LI MEDIUM & Early-to-late & DegreeOnly & \texttt{MEASURED\_\allowbreak{}FULL10} & 0.00568 / 253.30 & Measured, non-Pareto \\
LI MEDIUM & Early-to-late & NoEdge & \texttt{MEASURED\_\allowbreak{}FULL10} & 0.00515 / 243.66 & Measured, non-Pareto \\
LI MEDIUM & Early-to-late & RecentWindow & \texttt{MEASURED\_\allowbreak{}FULL10} & 0.00657 / 129.07 & Pareto \\
LI MEDIUM & Early-to-late & SAGE-edge h64 & \texttt{MEASURED\_\allowbreak{}FULL10} & 0.01030 / 249.24 & Pareto \\
LI MEDIUM & Early-to-late & ShuffledEdge & \texttt{MEASURED\_\allowbreak{}FULL10} & 0.00551 / 246.03 & Measured, non-Pareto \\
LI MEDIUM & Late-window & DegreeCap & \texttt{MEASURED\_\allowbreak{}FULL10} & 0.00256 / 251.89 & Measured, non-Pareto \\
LI MEDIUM & Late-window & DegreeOnly & \texttt{MEASURED\_\allowbreak{}FULL10} & 0.00778 / 303.13 & Measured, non-Pareto \\
LI MEDIUM & Late-window & NoEdge & \texttt{MEASURED\_\allowbreak{}FULL10} & 0.00600 / 294.95 & Measured, non-Pareto \\
LI MEDIUM & Late-window & RecentWindow & \texttt{MEASURED\_\allowbreak{}FULL10} & 0.00661 / 156.28 & Pareto \\
LI MEDIUM & Late-window & SAGE-edge h64 & \texttt{MEASURED\_\allowbreak{}FULL10} & 0.01105 / 300.67 & Pareto \\
LI MEDIUM & Late-window & ShuffledEdge & \texttt{MEASURED\_\allowbreak{}FULL10} & 0.00634 / 294.61 & Measured, non-Pareto \\
\end{longtable}
\noindent\footnotesize\emph{Scope and reading.} Pareto treatment is computed only within the same variant and protocol. Runtime is runner elapsed time, not a hardware-normalized complexity or energy measure. Blocked rows remain nonnumeric.\par
\endgroup


RecentWindow reduces aggregated Medium mean runtime from 275.19 to 141.65 seconds, while changing AUPRC by \(-0.01324\), AUROC by \(+0.02525\), and F1 by \(+0.00419\). Pareto flags are computed only within the same variant/protocol cell. Runtime is not normalized across hardware or implementations.

\paragraph{Sensitivity location.}
The complete context rows, including reference/candidate means, intervals, effect sizes, descriptive p-values, and source paths, remain at \path{results/tkde_rebuild/IBM_MATCHED_ABLATION_CONTEXT_EFFECTS.csv}. Their checksum and schema are reported in Table~\ref{tab:s-raw-artifacts}.
